\section{Discussion}
\label{sec:discussion}

\subsection{What the decomposition buys}

The literature on subject variability in physiological signal classifiers is
mostly a literature about remedies: add a subject identifier, adversarially
remove one, fine-tune on the new patient. Each remedy is expensive in a
different currency --- annotations, training time, or a labeled target
cohort. The decomposition of Section~\ref{sec:method} separates the patient
into a recording-wide offset and a stage-conditional signature, and the
experiments show that the two parts have different remedies. The offset is
manufactured by the normalization layer and can be removed by replacing its
statistics, with no optimizer and no label. The signature can be targeted only
by an operation that knows the stage, and the inability of any marginal method
to target it is visible in the within-stage probe.

The decomposition also changes what a subject-invariance claim should report.
An overall probe conflates the two components: a single vector subtraction
per recording drives it close to chance even when the conditional signature
is untouched, because the recording mean is the dominant direction. A paper
that reports only the overall probe would conclude that marginal adaptation
and SCSN are equivalent on HMC --- both reach a probe near $0.07$ --- when
the within-stage probe is $0.41$ under the former and $0.21$ under the
latter. The within-stage probe is the quantity that matches the object the
classifier uses: the relative position of two stage-conditional means.

\subsection{Why marginal adaptation is not enough}

Proposition~\ref{prop:marginal} gives the residual of any correction that
subtracts one vector per recording: $\mu_k+s_{\cR,k}+(\mu_{\mathrm{coh}}-\bar
\mu_{\cR})$. The first two terms say that the class-conditional means are
still those of the patient, not those of the cohort; the third says that the
correction adds an offset determined by the patient's stage mix. The second
term is why the within-stage probe survives. The third is why marginal
adaptation can \emph{lose} accuracy: on HMC, where recordings differ in stage
mix and the cohort is large, the post-hoc marginal arm is $2.3$ points worse
than the uncorrected model. The loss is not spread evenly over the night: on
HMC EEGNet the marginal arm gives up $0.135$ of REM recall and $0.107$ of
wake recall for $0.079$ of N2 recall, whereas SCSN raises N2 recall by
$0.040$ without a significant loss in any stage
(Section~\ref{sec:agreement}). An oracle marginal correction that knew the
recording's stage mix would remove the third term but not the second; the
signature is out of reach of any operation that uses one mean per recording.

SCSN removes both terms with the model's own posterior. The price is that the
posterior must be informative: the residual of the soft correction is the part
of the signature assigned to the wrong stage, bounded by the posterior error.
The cross-cohort experiment is the regime in which that price is not paid, and
the degradation there is the clearest evidence that the mechanism is the one
the algebra describes rather than an unspecified regularisation effect.

\subsection{What kind of claim this is}

The method can be read as an inference-time trick, and that reading would miss
what is being claimed. Nothing here changes the training procedure, the
architecture or the objective; the operation is arithmetic on a frozen
representation, and it is not presented as a new architecture or a new learning
objective. The claim is nevertheless about the model: that the
stage-conditional means of its penultimate representation are shifted by the
patient, that the shift is identifiable from the recording and the model's own
posterior, and that no operation using a single vector per recording can
remove it (Corollary~\ref{cor:no-marginal}). That is a statement about what the network
has learned, in the same sense as a probing result, and the paper makes it
falsifiable with a measurement it introduces: the within-stage probe falls
under SCSN and does not fall under any marginal correction, in every setting
tested. The estimator then follows from the statement through the conditional
expectation of the shift given the model's posterior
(Remark~\ref{rem:mmse}), and the closed form follows from the deployment
setting in which such a statement is useful: a checkpoint whose weights are
frozen and a night whose labels are unknown. A method that needed a second
pass, an optimizer or a gradient through the posterior would answer the same
analytical question and be a different engineering object. The object of the
contribution is the analysis and the measurement. The algorithm is what that
analysis leaves as the only available operation; its deployment properties are
what make the analysis usable at the bedside, and they are not where the
novelty lies.

\subsection{When the correction is free and when it is not}

On Sleep-EDF EEGNet, SCSN leaves accuracy where the uncorrected model and
layerwise AdaBN leave it; on HMC EEGNet and HMC SleepCNN it gains $1.9$ and
$1.4$ points over the uncorrected model and has the highest mean accuracy of
the arms, although on HMC SleepCNN the margin over layerwise AdaBN is not
significant ($p=0.084$).
The HMC gains survive the chance-corrected and class-balanced measures:
$\kappa$ rises by $0.023$ and $0.017$ and macro-$F_1$ by $0.018$ and $0.013$
on the two backbones, with no comparable change on Sleep-EDF.
On Sleep-EDF SleepCNN, layerwise AdaBN is $1.4$ points more accurate than
SCSN, a difference that is not significant. We report the ordering rather
than selecting the cohort on which the method wins. Two mechanisms plausibly
explain the ordering. Layerwise AdaBN is nonlinear and acts at every layer:
it removes part of the conditional signature as a side effect, lowering the
within-stage probe by $0.04$ to $0.09$ where it acts at all, and not at all
on Sleep-EDF SleepCNN, but it has no stage index and cannot target that
signature; SCSN removes $0.24$ to $0.35$ more. And the accuracy effect of
SCSN depends on the posterior: where the posterior is weak, the correction
approaches the marginal one, as the cross-cohort experiment shows. Neither
mechanism makes SCSN preferable on accuracy alone on Sleep-EDF; the case for
SCSN there rests on the identity it removes and on the fact that the same
operation is the only one available on a group-normalized checkpoint.

Tent is a strong baseline on Sleep-EDF and we do not claim to beat it on
accuracy. It is also a batch-normalization method: it optimizes the affine
parameters of batch-normalization layers. Applying it unchanged to a
group-normalized model is not defined, and a GroupNorm-specific entropy method
would require a separate design and tuning protocol. The comparison is
therefore reported for the batch-normalized checkpoints where both methods are
defined.

\subsection{The group-normalized setting}

Group normalization is often described as removing subject information
because it removes the batch. The experiments show that it does not remove
the patient: a linear probe identifies the patient from a group-normalized
representation well above chance on both cohorts and both backbones. What it
removes is the batch dependence of the normalization, not the
recording-dependent variation that the following layers map into the
penultimate space.

This is the setting in which the contribution is least ambiguous within the
family considered here. Standard AdaBN is undefined because there is no
test-time batch statistic, and standard Tent targets batch-normalization
affine parameters. The post-hoc marginal correction is defined and removes
the overall probe, but it moves the
within-stage probe by between a tenth and a third of what SCSN removes, and
in the wrong direction on both SleepCNN checkpoints, because a single vector
per recording carries no stage index. SCSN is the only stage-aware
statistics-replacement correction among those evaluated, and on the
group-normalized checkpoints it improves accuracy and the worst-covered
patient over the uncorrected model. A practitioner who chooses group
normalization for its batch independence therefore has a label-free
inference-time correction available without changing the architecture.

\subsection{Reading the reliability result correctly}

The conformal guarantee is marginal over patients, and the numbers make the
sense in which it is marginal concrete: every arm covers at about the nominal
rate on average, while the worst-covered patient of a fold is covered well
below it. The dispersion of per-patient coverage is the quantity that
separates a model whose guarantee is uniform from one whose guarantee is an
average.

The correction that removes the patient signature also reduces the
concentration of failures, but not uniformly across settings. On the
group-normalized checkpoints SCSN improves the worst-covered patient and the
dispersion in all four settings. On
the batch-normalized EEGNet checkpoints, layerwise AdaBN attains a better
worst-covered patient than SCSN ($0.866$ against $0.857$ on Sleep-EDF and
$0.698$ against $0.678$ on HMC) and a smaller dispersion on Sleep-EDF ($0.028$
against $0.036$), with the two dispersions comparable on HMC ($0.069$ against
$0.067$); on HMC SleepCNN the ordering reverses, with SCSN better on both
quantities ($0.664$ against $0.649$ and $0.066$ against $0.072$). We report
that rather than the comparison that favors the method. The safe statement is the mechanistic one: removing the
component of the representation that carries the patient reduces the
patient-dependence of the coverage, and the size of the reduction depends on
how much of that component the correction removes.

\subsection{The boundary of the method}

The cross-cohort experiment defines the regime in which a stage-conditional
correction should not be used. When a model trained on one centre is applied
to another, its posterior is weak, and the stage-conditional means estimated
from that posterior are dominated by the prior and by the recording offset.
SCSN then performs like the marginal correction and below layerwise AdaBN; the
oracle-label version, which uses the true stages to estimate the shifts,
recovers a large part of the gap, so the limitation is the posterior rather
than the estimator. This is a useful diagnostic: the method has a measurable
precondition, and the precondition can be checked on the unlabeled night
before the correction is applied, for example by the entropy or the
confidence of the uncorrected posterior.

\subsection{Limitations}

\paragraph{Cohort size}
Sleep-EDF has 20 subjects, so a fold holds out four of them and the probe is a
four-way problem. A four-way probe can only be so informative, and a
subject-level statistic computed over four subjects is noisy. This is the
reason the paper carries a second cohort of 151 subjects; the second cohort is
larger than most in the literature but comes from one centre, one montage and
one scoring team, and it is a cohort of patients referred for suspected sleep
disorders rather than of healthy volunteers.

\paragraph{Probe capacity}
The probe is a linear readout. A nonlinear probe would find more identity in
every arm, and the fractions reported here are lower bounds on what the
representation contains rather than statements about what it contains. The
comparison between arms is not affected by this, since the probe is the same
object in every arm, but the absolute level should be read as a floor.

\paragraph{Backbones}
Two convolutional backbones are used on both cohorts, chosen to differ in
depth, width and receptive field. Neither is a transformer, and the extent to
which normalization-driven identity matters in attention-based stagers is not
addressed here.

\paragraph{Posterior quality}
The correction is only as good as the posterior that indexes the stages. On
the within-cohort experiments the posterior is informative enough that the
correction is close to exact; across cohorts it is not, and the method
degrades to the marginal correction. A stager with a poorly calibrated head
would need the head recalibrated before the correction is applied.

\paragraph{Probability calibration}
Re-estimating the statistics, or correcting the representation, changes the
scale of the logits relative to the scale on which the posterior head was
fitted, and the expected calibration error moves with it: on Sleep-EDF EEGNet
from $0.030$ to $0.044$, with layerwise AdaBN at $0.043$ on the same runs, and
by at most $0.006$ on HMC. Conformal sets are
unaffected in coverage because their quantile is refitted on the calibration
split after the same intervention, but a user who reads the raw probabilities
of a deployed stager should recalibrate them after the correction.

\paragraph{Offline operation}
As with AdaBN, the statistics are estimated on the whole night. Clinical
staging is scored after the night, so the constraint is not binding in the
application we target, but a real-time monitor would have to estimate the
statistics on a prefix, which we have not tested.

\paragraph{GroupNorm-specific test-time adaptation}
Because group normalization has learnable affine parameters, an entropy-based
test-time objective could be designed for it. We did not implement that
baseline here: doing so would introduce a new optimizer, step-size schedule,
and stopping rule, and would no longer be the standard Tent method. The
group-normalized comparison therefore establishes applicability and internal
mechanism, not superiority to every possible GroupNorm-specific adaptation
objective.

\paragraph{Single channel}
Both cohorts are staged from a single EEG derivation. Multi-channel stagers
have more ways to encode a patient, and the balance between the marginal and
the conditional part of the identity may differ.

\paragraph{Statistical inference}
The Wilcoxon tests are exploratory and pairwise unadjusted. Seeds within a
fold reuse the same held-out subjects, so paired runs are not fully
independent; fold-level sign counts are therefore reported alongside the
run-level tests to reduce reliance on pseudo-replication. We make no
familywise-error claim. The consistency of effect directions across cohorts,
backbones, and normalization schemes, rather than any single $p$-value, is
the basis for the conclusions.

\subsection{Practical recipe}

For a stager that is already trained and validated, the procedure the results
support is short.
\begin{enumerate}
  \item Fit a linear probe from the penultimate representation to the patient,
  overall and within stage. If the overall probe is near chance, the marginal
  patient is not in the representation. If the within-stage probe is near
  chance, the conditional signature is not there either, and no
  stage-conditional correction is needed.
  \item If the within-stage probe is high, apply SCSN on each recording before
  predicting it. On a batch-normalized model, compare with layerwise AdaBN; on
  a group-normalized model, SCSN is the only stage-conditional option.
  \item Check the uncorrected posterior before applying the correction. If it
  is close to uniform, the stage-conditional means are not identifiable and
  the correction should be reduced to the marginal one.
  \item Recalibrate the conformal score on recordings processed the same way,
  and report per-patient coverage dispersion beside average coverage.
\end{enumerate}
The first step is a diagnostic that costs one logistic regression; the second
adds a few matrix products to the forward pass the stager performs anyway; the
third is a scalar check on the posterior; the fourth is bookkeeping. None of them requires a labeled target cohort, and
all of them can be applied to a model whose weights will never be touched
again.
